\section{Graph-conditional revalidation}
\label{sec:method}
We distinguish a graph's predictive credit from the value of purchasing another
measurement. For a question $x$, generation stream $s$, and communication graph
$G$, let $Y(x,G,s)\in\{0,1\}$ be exact-match correctness. We use deletion gain
\begin{equation}
 d_e(x,G,s)=Y(x,G\setminus e,s)-Y(x,G,s).
\end{equation}
Positive deletion gain favors removing the edge. This sign is the opposite of
an edge's contribution to correctness; the convention must not be reversed
when comparing historical credit-prediction results. A graph change can alter
$d_e$ without changing the edge $e$ itself. Neither a nonzero interaction nor a
detectable credit change establishes that buying feedback is worthwhile.

\paragraph{Controlled graph change.}
The four nodes are $U,V,R,M$, with $M$ producing the answer. Nodes $U,V,M$ use
Qwen2.5-3B-Instruct and $R$ uses Mistral-7B-Instruct-v0.3. Initially,
\begin{align}
 E_0=\{&U\!\to\!R,V\!\to\!R,R\!\to\!M,\nonumber\\
       &U\!\to\!M,V\!\to\!M\}.
\end{align}
A predetermined exogenous shock removes either $U\to R$ or $V\to R$.
The action menu keeps the resulting graph or deletes one direct edge,
$e_U=U\to M$ or $e_V=V\to M$. Every executed node receives all ten raw
documents in the distractor input. The operator executes exactly the terminal
node and its ancestors. Thus deleting the affected direct edge also removes an
unused source computation. This execution rule is shared by every method;
dead-node elimination is not claimed as our contribution.

\paragraph{Information separation.}
Each episode has eight initialization questions $I$, eight feedback questions
$C$, and 64 future service questions $S$. Initialization provides paired gains
under $G_0$. A policy receives an outcome from $C$ only after selecting and
paying for its paired counterfactual. All terminal decisions are fixed before
the evaluator uses $S$. The experiment measures a full counterfactual archive
once, but the policy interface conceals all unpurchased feedback. This is
controlled cache replay, not a deployed online system. Benchmark reference
answers serve as an idealized feedback oracle.

\paragraph{A fixed working predictor.}
For each editable edge we use a ternary model over $d_e\in\{-1,0,1\}$ with
Dirichlet pseudocount $1/3$ per outcome, updated by all eight initialization
observations. Purchased current-graph observations add one count. Old counts
are retained without graph-dependent discounting. For counts $\alpha_{e,j}$,
\begin{equation}
 \mu_e=\frac{\alpha_{e,+1}-\alpha_{e,-1}}{\sum_j\alpha_{e,j}}.
\end{equation}
These are working estimates, not calibrated guarantees under graph shift.
The model is deliberately simple and identical across acquisition strategies.
The experiment does not propose a new predictor architecture.

\paragraph{Graph-dependent service prices.}
Let $c_0,c_U,c_V$ be mean token costs of the initial keep and deletion graphs,
estimated only from $I$. Under the changed graph, the source-call price is
subtracted from the affected deletion graph's forecast because that node is
no longer an ancestor of $M$. No future service cost or outcome is used in
this forecast. With fixed token weight $\lambda=0.005/1000$, the action values
relative to keeping are
\begin{equation}
 u_0=0,\qquad u_e=\mu_e+\lambda(c_0-c_e).
\end{equation}
The maximum selects the terminal action; ties favor keep, then $U$, then $V$.
The weight is a declared research tradeoff, not a monetary conversion.

\paragraph{Acquisition and stopping.}
Write $V(\alpha)=\max_a u_a(\alpha)$ and let
\begin{equation}
 \mathrm{EVSI}_e=\sum_{j=-1}^{1}
 \frac{\alpha_{e,j}}{\sum_k\alpha_{e,k}}
 V(\alpha+\mathbf{1}_{e,j})-V(\alpha).
\end{equation}
The proposed nominee ranks an available edge's next feedback question by
$64\,\mathrm{EVSI}_e/\widehat m_{e,q}$, where $\widehat m_{e,q}$ is its
forecast incremental token price. The price includes the current base graph
and counterfactual final call; if that question's base execution has already
been purchased, the base forecast is omitted. The forecasts use only $I$ and
the paid cache state. Actual charged cost is determined by exact request
fingerprints after purchase. The policy stops if its selected measurement has
$64\,\mathrm{EVSI}_e\leq\lambda\widehat m_{e,q}$, or after eight purchases.
EVSI, cost normalization, and stopping are established design principles
\citep{diva}; the tested hypothesis concerns their incremental value with
graph-dependent execution and feedback prices.

\paragraph{Ablations and free controls.}
The fixed-price EVSI control uses the same predictor and current service
prices but always forecasts the full base-plus-counterfactual acquisition
price. Other ablations remove cost normalization from the ranking or use
old-graph service prices. Random, periodic, maximum-uncertainty, and
changed-edge-first strategies each purchase eight pairs. Stale credit makes
no purchase; protect-changed forbids deleting the affected direct edge;
keep always keeps. All these maintenance methods pay initialization.
Three additional static controls keep, delete the changed edge, or delete
the other edge without paying initialization. They prevent free structural
information from being mistaken for a benefit of paid revalidation.

\paragraph{System baselines.}
Single Qwen sees the full raw input. Qwen voting uses three full-context
answer generations with normalized plurality and fixed tie order, returning
the selected complete raw answer. Independent-three aggregation uses two
Qwen sources and one Mistral source, all seeing the full input, followed by
a full-context Qwen aggregator that reads their notes (four calls in total).
The suffix Q in this baseline identifies its final model, not three
homogeneous Qwen sources. Its prompts, model allocation, and seed rules were
checked against the preceding engineering baseline.
